A região do Estreito de Ormuz é um dos pontos mais críticos do comércio global, exigindo monitoramento constante e capacidade de resposta rápida a incidentes marítimos. Tradicionalmente, sistemas de monitoramento operam sob modelos centralizados, onde uma falha no servidor central pode paralisar toda a operação de vigilância e resgate. O projeto HormuzNet surge para endereçar essa vulnerabilidade através de uma arquitetura de sistemas distribuídos descentralizada.

O objetivo principal deste trabalho é descrever a implementação de uma malha autônoma e resiliente, onde múltiplos \textit{brokers} cooperam para processar dados de sensores ambientais e coordenar frotas de drones autônomos. A solução foca na eliminação de pontos únicos de falha e na garantia de consistência eventual através de algoritmos clássicos de sistemas distribuídos, como os Relógios de Lamport.

O sistema é composto por três tipos de entidades: sensores, que publicam dados via UDP Multicast; \textit{brokers}, que formam a malha de controle P2P; e drones autônomos, que executam missões de resposta. Este documento detalha as decisões de projeto, a implementação da concorrência em linguagem Go, os protocolos de rede utilizados e os resultados obtidos na validação da resiliência do sistema.
